\documentclass[12pt,a4paper]{article}
% Packages
\usepackage[utf8]{inputenc}
\usepackage[T1]{fontenc}
\usepackage[french]{babel}
\usepackage{amsmath, amssymb}
\usepackage{geometry}
\usepackage{graphicx}
\usepackage{xcolor}

\geometry{margin=1.5cm, top=0.5cm}

\begin{document}
	
	% ================= EN-TÊTE =================
	\begin{tabular}{p{3cm}p{11cm}p{3cm}}
		&\begin{center}
			\includegraphics[width=10cm,height=0.3cm]{../../Pictures/IMG-20260408-WA0005}
		\end{center}&\\
		\vspace{0.3cm}
	\includegraphics[width=4cm,height=4cm]{../../Pictures/IMG-20260408-WA0006}	& 
		\begin{center}
			\textbf{\large République du Bénin}\\
			\hspace{-0.5cm}\textbf{\large Ministère de l'Enseignement }%\\[0.1cm]
			\textbf{\large Supérieur et de la Recherche  }%\\[0.1cm]
			\textbf{\large Scientifique (MESRS)}\\ \vspace{0.22cm}
			
			\textbf{ \large	Université Nationale des Sciences, Technologies, Ingénierie et Mathématiques (UNSTIM) d'Abomey} \\ [0.1cm]
			
			\textbf{ÉCOLE NORMALE SUPÉRIEURE DE NATITINGOU}
		\end{center}
		& \vspace{0.5cm} \includegraphics[width=3cm,height=3cm]{../../Pictures/IMG-20260408-WA0004}
	\end{tabular}
	\noindent \textcolor{cyan}{\rule{\textwidth}{2.5pt}}
	% ================= PARTIE I =================
	\section*{I. MATHÉMATIQUES (7 pts)}
	\noindent \textcolor{blue}{\rule{\textwidth}{2.5pt}}
	
	Soit $f$ une fonction continue périodique. On considère ses coefficients :
	\[c_n(f) = \frac{1}{2\pi} \int_0^{2\pi} f(t)e^{-int}\,dt
	\]
	
	\begin{enumerate}
		\item Prouver que
		\[	\sum_{n=-\infty}^{+\infty} |c_n(f)|^2 
		= \frac{1}{2\pi} \int_0^{2\pi} |f(t)|^2 \, dt.	\]
	\end{enumerate}
	
	\vspace{0.5cm}
	
	% ================= PARTIE II =================
	\section*{II. PHYSIQUE (7 pts)}
	\noindent \textcolor{red}{\rule{\textwidth}{2.5pt}}
	
	Démontrer l’invariant de la pseudo-norme du quadrivecteur impulsion :
	\[	\left(\frac{E}{c}\right)^2 - \overrightarrow{p }^2 = m^2 c^2	\]
	
	\vspace{0.5cm}
	% ================= PARTIE III =================
	\section*{III. CHIMIE (6 pts)}
	\noindent \textcolor{green}{\rule{\textwidth}{2.5pt}}
	
	Calculer l’enthalpie libre de réaction $\Delta_r G$ en fonction du quotient de réaction $Q$ :
	\[	\Delta_r G = \Delta_r G^\circ + RT \ln(Q)	\]
	
	\vspace{1cm}
	
	% ================= TEXTE SURIGNÉ =================
	\noindent
	\colorbox{yellow}{
		\parbox{\textwidth}{
			En Afrique subsaharienne, la transformation digitale en santé se déploie dans un contexte singulier. Les systèmes de santé y sont caractérisés par une insuffisance chronique de ressources financières, humaines et matérielles(Nobyonga-Orem et Taslobya,2016). Les infrastructures numériques, bien que croissantes avec l’essor de la téléphonie mobile et d’Internet, restent inégales et parfois précaires. Les coupures d’électricité fréquentes, l’instabilité des réseaux et la faible capacité technique du personnel constituent des contraintes majeures(Dovlo,2018) . Cependant, ces fragilités coexistent avec un réel potentiel.
		}
	}
	
	\vspace{1cm}
	
	\begin{flushright}
		\textbf{Page 1 / 1}
	\end{flushright}
	

	
\end{document}